\section*{Chapter \ref{ch:g11:gene-expression} --- Gene Expression}

\begin{solution}{exo:g11:gene-expression:1}
The central flow of genetic information for expression is \omterm{def:g11:dna-replication:helix}{DNA} $\to$ \omterm{def:g10:molecules-of-life:nucleic}{RNA} $\to$
\omterm{def:g10:molecules-of-life:proteins}{protein} (\omterm{def:g11:gene-expression:transcription}{transcription} then \omterm{def:g11:gene-expression:translation}{translation}). \omterm{def:g11:dna-replication:semi}{Replication} is \omterm{def:g11:dna-replication:helix}{DNA} $\to$ \omterm{def:g11:dna-replication:helix}{DNA}, copying
the archive for daughter \omterm{def:g10:the-cell:cell}{cells}; it is not part of expressing a \omterm{def:g11:gene-expression:gene}{gene} into
product, but it uses the same base-pairing logic.
\end{solution}

\begin{solution}{exo:g11:gene-expression:2}
The mRNA \omterm{def:g10:molecules-of-life:proteins}{sequence} \omterm{prop:g11:dna-replication:base-pairs}{complementary} to the given template, written $5'\to 3'$, is
\texttt{5'-AUG UUU GGC UAA-3'} (spaces optional for readability). \omterm{def:g10:molecules-of-life:nucleic}{RNA} uses U
in place of T.
\end{solution}

\begin{solution}{exo:g11:gene-expression:3}
Among the standard bases, uracil (U) appears in \omterm{def:g10:molecules-of-life:nucleic}{RNA} and thymine (T) in \omterm{def:g11:dna-replication:helix}{DNA}.
Both pair with adenine; the sugar and base sets distinguish the two \omterm{def:g10:molecules-of-life:polymer}{polymers}.
\end{solution}

\begin{solution}{exo:g11:gene-expression:4}
The \omterm{def:g11:gene-expression:strands}{template strand} is read by \omterm{def:g11:gene-expression:transcription}{RNA polymerase} and is \omterm{prop:g11:dna-replication:base-pairs}{complementary} to the \omterm{def:g10:molecules-of-life:nucleic}{RNA}.
The coding (non-template) strand has the same \omterm{def:g10:molecules-of-life:proteins}{sequence} as the \omterm{def:g10:molecules-of-life:nucleic}{RNA} except that
\omterm{def:g11:dna-replication:helix}{DNA} has T where \omterm{def:g10:molecules-of-life:nucleic}{RNA} has U, written in the $5'\to 3'$ direction; it is not the
strand used as the polymerase template.
\end{solution}

\begin{solution}{exo:g11:gene-expression:5}
\omterm{def:g11:gene-expression:trna}{Anticodon} pairs \omterm{def:g11:dna-replication:antiparallel}{antiparallel} with the \omterm{def:g11:gene-expression:codon}{codon}. \omterm{def:g11:gene-expression:codon}{Codon} \texttt{5'-GCA-3'} pairs
with \omterm{def:g11:gene-expression:trna}{anticodon} \texttt{3'-CGU-5'}, which is \texttt{5'-UGC-3'} when rewritten
$5'\to 3'$. Base pairing at the \omterm{def:g11:gene-expression:translation}{ribosome} enforces the \omterm{def:g11:gene-expression:codon}{genetic code}.
\end{solution}

\begin{solution}{exo:g11:gene-expression:6}
The code is degenerate: most \omterm{def:g10:molecules-of-life:proteins}{amino acids} have several \omterm{def:g11:gene-expression:codon}{codons}. A base change may
yield a synonymous \omterm{def:g11:gene-expression:codon}{codon} and leave the \omterm{def:g10:molecules-of-life:proteins}{amino acid} unchanged --- a silent
mutation at the \omterm{def:g10:molecules-of-life:proteins}{protein} \omterm{def:g10:molecules-of-life:proteins}{sequence} level (expression level might still change
slightly).
\end{solution}

\begin{solution}{exo:g11:gene-expression:7}
mRNA: \texttt{5'-AUG GGA UAA-3'}. \omterm{def:g11:gene-expression:translation}{Translation} begins at AUG (Met), adds Gly
from GGA, then hits UAA stop --- a dipeptide Met--Gly, not a longer chain.
\end{solution}

\begin{solution}{exo:g11:gene-expression:8}
\omterm{prop:g11:gene-expression:tx-stages}{Initiation}: the \omterm{def:g11:gene-expression:translation}{ribosome} assembles at the \omterm{prop:g11:gene-expression:code-props}{start codon} with initiator \omterm{def:g11:gene-expression:translation}{tRNA}.
\omterm{prop:g11:gene-expression:tx-stages}{Elongation}: \omterm{def:g10:molecules-of-life:proteins}{peptide bonds} form as the \omterm{def:g11:gene-expression:translation}{ribosome} advances \omterm{def:g11:gene-expression:codon}{codon} by \omterm{def:g11:gene-expression:codon}{codon} with
\omterm{def:g11:gene-expression:translation}{aminoacyl-tRNAs}. \omterm{prop:g11:gene-expression:tx-stages}{Termination}: a \omterm{prop:g11:gene-expression:code-props}{stop codon} recruits release factors and frees
the polypeptide.
\end{solution}

\begin{solution}{exo:g11:gene-expression:9}
Insertion of one base shifts the reading frame: almost every \omterm{def:g11:gene-expression:codon}{codon} downstream
is wrong, usually producing a garbled and/or truncated \omterm{def:g10:molecules-of-life:proteins}{protein}. A synonymous
substitution changes \omterm{def:g11:dna-replication:helix}{DNA} but not the amino-acid \omterm{def:g10:molecules-of-life:proteins}{sequence}, so \omterm{def:g10:molecules-of-life:proteins}{primary structure}
is unchanged.
\end{solution}

\begin{solution}{exo:g11:gene-expression:10}
The primary eukaryotic transcript still contains \omterm{prop:g11:gene-expression:tx-stages}{introns} and lacks mature ends.
Processing includes $5'$ capping, $3'$ polyadenylation, and \omterm{prop:g11:gene-expression:tx-stages}{splicing} (intron
removal and exon joining) before export and \omterm{def:g11:gene-expression:translation}{translation}.
\end{solution}

\begin{solution}{pb:g11:gene-expression:1}
\textbf{1.} Template: \texttt{3'-TAC GGA CTT AAG ACC ATC CCT-5'}.

\textbf{2.} mRNA: \texttt{5'-AUG CCU GAA UUC UGG UAG GGA-3'}.

\textbf{3.} Met--Pro--Glu--Phe--Trp, then stop (\texttt{UAG}). The trailing
\texttt{GGA} is not translated.

\textbf{4.} Coding \texttt{TGA} $\to$ mRNA \texttt{UGA}, a \omterm{prop:g11:gene-expression:code-props}{stop codon}: the
\omterm{def:g10:molecules-of-life:proteins}{protein} ends after Phe (truncated; Trp lost).

\textbf{5.} Deleting one base shifts the frame from that point: wrong \omterm{def:g10:molecules-of-life:proteins}{amino
acids} and likely a premature or delayed stop --- a frameshift.

\textbf{6.} Bacterium: \omterm{def:g11:gene-expression:transcription}{transcription} and \omterm{def:g11:gene-expression:translation}{translation} both in the cytoplasm (no
nucleus). Human \omterm{def:g10:the-cell:cell}{cell}: \omterm{def:g11:gene-expression:transcription}{transcription} and processing in the nucleus; \omterm{def:g11:gene-expression:translation}{translation}
on cytosolic \omterm{def:g11:gene-expression:translation}{ribosomes} (or rough ER).

\textbf{7.} No \omterm{def:g10:the-cell:nucleus}{nuclear envelope} separates \omterm{def:g11:dna-replication:helix}{DNA} from \omterm{def:g11:gene-expression:translation}{ribosomes} in bacteria, so
\omterm{def:g11:gene-expression:translation}{ribosomes} can bind nascent mRNA. In human \omterm{def:g10:the-cell:cell}{cells} mRNA must be processed and
exported from the nucleus before cytosolic \omterm{def:g11:gene-expression:translation}{translation}.
\end{solution}
